\textbf{Problème 1.15 (Dujella, 2026).} Existe-t-il un quadruplet diophantien fort, c'est-à-dire un ensemble de quatre rationnels non nuls $\{a_1,a_2,a_3,a_4\}$ tels que $a_ia_j+1$ soit un carré parfait pour tout $1 \leqslant i,j \leqslant 4$ ?

\medskip

\textit{Remarque.} Il existe des exemples de quadruplets diophantiens « presque forts », où toutes les conditions d'un quadruplet fort sont vérifiées sauf que $a_3a_4+1$ n'est pas un carré, par exemple $\{\frac{140}{51}, \frac{2223}{3046}, \frac{278817}{33856}, \frac{3182740}{17661}\}$ (Dujella-Petričević, Experiment. Math. \textbf{17} (2008), 83--89), mais on ne sait pas s'il en existe une infinité.
